% Counterexample.tex — companion note to Counterexample.lean
% (formal verification source: Proofs.Erdos20.Counterexample in
%  github.com/thatnealpatel/proofs).
\documentclass[11pt,a4paper]{article}

\usepackage[T1]{fontenc}
\usepackage{amsmath,amssymb,amsthm}

\newtheorem{theorem}{Theorem}
\newtheorem{corollary}{Corollary}
\newtheorem{definition}{Definition}
\theoremstyle{remark}
\newtheorem{remark}{Remark}

\newcommand{\cC}{\mathcal{C}}
\newcommand{\cF}{\mathcal{F}}
\newcommand{\cG}{\mathcal{G}}
\newcommand{\cS}{\mathcal{S}}

\title{Frankl Shifts Can Create Sunflowers:\\
  a Machine-Checked Counterexample to a Proposed Bridge\\
  for the Sunflower Conjecture}
\author{Neal Patel\thanks{\texttt{neal@patel.codes}. The Lean development
  and this note were produced with AI assistance; see the AI disclosure in
  the repository (\texttt{Proofs/README}).}}
\date{June 10, 2026}

\begin{document}
\maketitle

\begin{abstract}
Version~1 of Mishra, \emph{Erd\H{o}s Rado sunflower (conjecture)
theorem}%
\footnote{The v1 title is ``Erd\H{o}s Rado Sunflower (Conjecture) Theorem'';
  retitled in v2 to ``Erd\H{o}s Rado Sunflower Theorem for Shifted Families.''}
(arXiv:2606.02667v1), proposed reducing the Erd\H{o}s--Rado
sunflower conjecture to the case of shifted families via the bound
$\tau(\cS(\cF))\le3\,\tau(\cF)^{2}$ (its Lemma~3), where $\cS(\cF)$ is
the endpoint of the canonical Frankl-shift sweep and $\tau$ denotes the
largest size of a sunflower contained in a family. We exhibit a
$4$-uniform family $\cF$ on $20$ points with $17$ members and
$\tau(\cF)=2$ whose canonical sweep endpoint is a $17$-petal star, so
that $\tau(\cS(\cF))=17>12=3\,\tau(\cF)^{2}$. The same witness refutes
the natural repaired bound $3\tau^{2}+\tau+1=15$. All claims are
machine-checked in Lean~4: the chain endpoint, its stability,
$\tau(\cF)=2$ (by exhaustive scan), and $\tau(\cS(\cF))=17$. The
mechanism is that shifting funnels low-$\tau$ families into stars on
the lexicographically smallest $k-1$ elements; only empty-kernel
sunflowers (matchings) are guaranteed to transport backward through
shift chains.
\end{abstract}

\section{The proposed bridge}

A \emph{sunflower} of size $s$ with kernel $Y$ is a collection of $s$
distinct sets whose pairwise intersections all equal $Y$ and whose
petals (members minus $Y$) are non-empty; $\tau(\cF)$ is the largest
size of a sunflower contained in $\cF$, and $\nu(\cF)\le\tau(\cF)$ is
the matching number. For $i<j$, the Frankl compression $\cC_{ij}$
\cite{Fr87} replaces $j$ by $i$ in each member $S$ of a family $\cF$
(when $i\notin S$, $j\in S$) unless the replacement already lies in
$\cF$, in which case $S$ is kept; $\cC_{ij}$ is injective on $\cF$, so
$|\cC_{ij}(\cF)|=|\cF|$. The \emph{canonical sweep} applies all
$\binom{n}{2}$ compressions in lexicographic order, and its endpoint
is denoted $\cS(\cF)$; a family is \emph{shifted} when every
$\cC_{ij}$ with $i<j$ fixes it.

Version~1 of \cite{Mi26} claimed the full Erd\H{o}s--Rado conjecture
\cite{ER60} by combining a (correct, see the companion note
\texttt{ShiftedSunflower.tex}) sunflower theorem for shifted families,
$f'(k,s)\le s^{2s-2}2^{k}$, with a bridge from general families to
shifted ones:

\begin{quote}
\textbf{Lemma~3 of \cite{Mi26}v1.}\quad
$\tau(\cS(\cF))\;\le\;3\,\tau(\cF)^{2}$.
\end{quote}

\noindent Together these would give $f(k,s)\le(3s^{2})^{6s^{2}-2}\,2^{k}$
(v1's Theorem~2). This note records a machine-checked proof that the
bridge's \emph{statement} --- not merely its proof --- is false.
Version~2 of \cite{Mi26} withdraws Lemma~3 and the unconditional claim
(both are commented out in the v2 source) and retains the shifted-case
theorem.

\section{The witness}

Elements are written $0$-indexed over $\{0,\dots,19\}$, exactly as in
the Lean development (\texttt{Fin 20}); add $1$ to every element for
the $[n]$ convention of \cite{Mi26}. Let $\cF$ be the following
$4$-uniform family with $17$ members (\texttt{familyB} in the Lean
source):
\[
\begin{array}{lllll}
\{0,1,2,15\} & \{0,1,3,15\} & \{0,1,4,5\} & \{0,4,14,16\} & \{0,4,14,18\}\\
\{0,8,9,15\} & \{0,8,10,15\} & \{1,2,3,4\} & \{1,2,3,7\} & \{1,4,5,8\}\\
\{1,5,6,12\} & \{2,9,11,12\} & \{2,3,12,13\} & \{2,3,6,12\} & \{3,11,12,17\}\\
\{3,11,14,15\} & \{4,14,15,19\}
\end{array}
\]

\begin{theorem}[machine-checked]\label{thm:witness}
\leavevmode
\begin{enumerate}
\item $\tau(\cF)=2$: the pair $\{0,1,2,15\},\{0,1,3,15\}$ is a
$2$-sunflower with kernel $\{0,1,15\}$, and no three members of $\cF$
form a sunflower (exhaustive scan over all
$\binom{17}{3}=680$ triples).
\item The canonical sweep --- all $\cC_{ij}$ for
$0\le i<j\le19$ in lexicographic order ($i$ ascending, $j$ ascending
within $i$; $190$ compressions, the membership guard evaluated against
the current family at each step) --- sends $\cF$ to the full star
\[
  \cS(\cF)\;=\;\bigl\{\{0,1,2,x\}:3\le x\le19\bigr\}.
\]
\item $\cS(\cF)$ is stable under every $\cC_{ij}$ with $i<j$.
\item $\tau(\cS(\cF))=17$: the $17$ members pairwise intersect in
exactly $\{0,1,2\}$ with distinct singleton petals, so the star is
itself a $17$-sunflower, and $\tau\le|\cS(\cF)|=17$.
\end{enumerate}
\end{theorem}

\begin{corollary}\label{cor:false}
$\tau(\cS(\cF))=17>15=3\,\tau(\cF)^{2}+\tau(\cF)+1>12=3\,\tau(\cF)^{2}$.
Lemma~3 of \cite{Mi26}v1 is false, and so is its natural repair
$\tau(\cS(\cF))\le3\tau^{2}+\tau+1$.
\end{corollary}

The Lean statements are
\texttt{shifting\_creates\_stars} --- there exist families
$\cF,\cG$ over \texttt{Fin 20} with $\cG$ a full shift of $\cF$,
$\tau(\cF)=2$, and $\tau(\cG)>15$ --- and
\texttt{mishra\_v1\_lemma3\_false} --- the universally quantified bound
$\tau(\cG)\le3\tau(\cF)^{2}+\tau(\cF)+1$ over full-shift pairs is
refuted. Here ``$\cG$ is a full shift of $\cF$''
(\texttt{IsFullShiftOf}) means: $\cG$ is reachable from $\cF$ by a
chain of $(i<j)$-compressions and is stable under all of them. By
Theorem~\ref{thm:witness}(2,3) the canonical sweep endpoint is such a
$\cG$, so the refutation applies verbatim to the $\cS(\cdot)$ of
\cite{Mi26}.

\section{Mechanism and scope}

\paragraph{Shifting creates stars.}
The sweep collapses all $17$ members onto the common core $\{0,1,2\}$
--- the lexicographically smallest $k-1$ elements --- manufacturing a
maximal star. Sunflowers with non-empty kernel are thereby created
from nothing: $\cF$ has no sunflower of size $3$, while $\cS(\cF)$
contains one of size $17$.

\paragraph{Matchings are immune.}
By contrast, empty-kernel sunflowers cannot be created: a matching of
size $t$ in any chain endpoint pulls back to a matching of size $t$ in
the source. This is formalized
(\texttt{matchingNumber\_shift\_le},
\texttt{matchingNumber\_reachable} in
\texttt{Proofs.Erdos20.Sunflower}; the single-step statement is
\cite[Lemma~1(iii)]{Mi26}, $\nu(\cC_{ij}(\cF))\le\nu(\cF)$). The
failure of the bridge is therefore confined exactly to the
non-empty-kernel case: here $\nu(\cF)\le\tau(\cF)=2$ bounds
$\nu(\cS(\cF))$, while $\tau(\cS(\cF))=17$ --- the chain-stable
invariant $\nu$ cannot see the blow-up.

\paragraph{Further witnesses.}
Three more witnesses were certified by two independently written
verification programs (they are not formalized in Lean): at
$(k,n,|\cF|)=(4,16,21)$ the endpoint contains the complete star on the
first three elements, of size $13$, refuting the $3\tau^{2}=12$ bound
but not the repaired $15$; at $(5,22,30)$ and $(6,24,43)$ the
endpoints contain stars of size $17$, refuting both. In all cases
$\tau(\cF)=2$ by exhaustive scan. At $k=4$ the two certified endpoints
carry the complete star of size exactly $n-3$ ($13$ at $n=16$, $17$ at
$n=20$); if that growth persists in $n$ --- search experiments suggest
it does, but no larger instance has been certified --- then no bound
of the form $\tau(\cS(\cF))\le g(\tau(\cF),k)$ can hold with $g$
independent of the ground-set size.

\paragraph{What the witnesses do not exclude.}
All witnesses are small ($|\cF|\le43$), far below the shifted-case
threshold $s^{2s-2}2^{k}$; a bridge hypothesis restricted to families
\emph{above} the sunflower-free size threshold is untested. Bridges
through invariants other than $\tau$ are likewise untouched, with one
caveat: on the shifted side, bounds $\tau(\cG)\le g(\nu(\cG))$ for
shifted $\cG$ fail trivially, since full stars are shifted families
with $\nu=1$ and $\tau$ as large as the ground set allows.

\section{The Lean certificate}

The module \texttt{Proofs.Erdos20.Counterexample} (next to this note)
imports the shift theory of \texttt{Proofs.Erdos20.Sunflower}
(\texttt{franklShift}, \texttt{IsFullyCompressed},
\texttt{IsFullShiftOf}, \texttt{sunflowerNumber}) and proceeds as
follows.

\begin{itemize}
\item \emph{Computable twin.} \texttt{franklShift} is marked
noncomputable; a twin \texttt{franklShiftC} with an identical body is
definitionally equal (proved by \texttt{rfl}), so evaluated facts about
the twin convert to the original for free.
\item \emph{The chain as data.} The $190$ ordered pairs are a literal
list; \texttt{applyChain} folds the compressions left-to-right, and a
structural induction (\texttt{reachable\_applyChain}) turns any such
list into a reachability certificate. \texttt{native\_decide}
evaluates the ground facts: the endpoint equality
$\texttt{applyChain sweep familyB}=\texttt{starB}$, the stability of
\texttt{starB} under all $(i<j)$-compressions, and
$|\texttt{starB}|=17$.
\item \emph{Deciding $\tau(\cF)=2$.} The predicate ``$\cF$ has an
$s$-sunflower'' quantifies existentially over arbitrary subfamilies
and kernels and is not directly decidable at feasible cost; but any
$3$-sunflower yields a bounded witness: a $3$-element subfamily
together with a kernel contained in any one of its members. The
bounded form is refuted by \texttt{native\_decide} ($680$ triples,
at most $2^{4}=16$ candidate kernels each); the $2$-sunflower is
exhibited literally; attainment of the supremum closes
$\tau(\cF)=2$. A practical note: the bounded quantifiers must be
decided by \texttt{Finset} traversal --- Lean's default instance
search prefers an instance that enumerates all $2^{20}$ subsets of the
ground set per quantifier; pinning the traversal instances at high
local priority makes the certificate evaluate in seconds.
\item \emph{Trust base.} \texttt{\#print axioms} reports
\texttt{propext}, \texttt{Classical.choice}, \texttt{Quot.sound},
plus one trust axiom per \texttt{native\_decide} evaluation (the
Lean~4.30 form of \texttt{Lean.ofReduceBool}: the compiler is trusted
for eight ground, finitely checkable facts). There is no
\texttt{sorry}. All reasoning around the ground facts --- the $\tau$
algebra, the reachability certificate, the final inequalities --- is
ordinary kernel-checked Lean.
\item \emph{Provenance.} The witness was found by stochastic search
and certificate-verified by two independently written programs
(disjoint code, chain semantics matched to the Lean definition of
\texttt{franklShift}) before formalization; the Lean module
re-verifies every certified claim inside the proof assistant.
\end{itemize}

\begin{center}
\begin{tabular}{ll}
\hline
Claim & Lean identifier \\
\hline
the witness family and the star &
  \texttt{familyB}, \texttt{starB} \\
Theorem~\ref{thm:witness}(1) &
  \texttt{familyB\_no\_3sunflower}, \texttt{familyB\_hasSunflower\_2},
  \texttt{familyB\_tau} \\
Theorem~\ref{thm:witness}(2) &
  \texttt{sweep}, \texttt{chain\_endpoint},
  \texttt{familyB\_reaches\_starB} \\
Theorem~\ref{thm:witness}(3) &
  \texttt{starB\_compressed} \\
Theorem~\ref{thm:witness}(4) &
  \texttt{starB\_hasSunflower\_17}, \texttt{starB\_tau} \\
Corollary~\ref{cor:false} &
  \texttt{shifting\_creates\_stars}, \texttt{mishra\_v1\_lemma3\_false} \\
\hline
\end{tabular}
\end{center}

\begin{thebibliography}{Mi26}

\bibitem[ER60]{ER60}
P.~Erd\H{o}s and R.~Rado,
\emph{Intersection theorems for systems of sets},
J.\ London Math.\ Soc.\ \textbf{35} (1960), 85--90.

\bibitem[Fr87]{Fr87}
P.~Frankl,
\emph{The shifting technique in extremal set theory},
in: Surveys in Combinatorics 1987, London Math.\ Soc.\ Lecture Note
Ser.\ \textbf{123}, Cambridge Univ.\ Press, 1987, 81--110.

\bibitem[Mi26]{Mi26}
T.~K.~Mishra,
\emph{Erd\H{o}s Rado sunflower (conjecture) theorem},
arXiv:2606.02667 (v1: 2026-06-01, titled ``Erd\H{o}s Rado Sunflower
(Conjecture) Theorem''; v2: 2026-06-08, retitled ``Erd\H{o}s Rado
Sunflower Theorem for Shifted Families'').

\bibitem[MU21]{MU21}
L.~de~Moura and S.~Ullrich,
\emph{The Lean 4 theorem prover and programming language},
in: Automated Deduction --- CADE 28, LNCS \textbf{12699}, Springer,
2021, 625--635.

\end{thebibliography}

\end{document}
